\begin{abstract}
Coherent QCA was proposed as a route to quantum computation almost
twenty years ago~\cite{wu2005coherent,xu2008universal}, but a full
fault-tolerant stack on this platform has not been demonstrated.
We build one, and extend the same lattice to two classes of
programmable quantum simulation.

(i) We confirm the native QCA coupling is locally a CNOT and
prepare a Bell state at fidelity $0.9975$ and a GHZ chain at
$0.9999$.  (ii) We show that the four-phase clock is bounded by
the Landau--Zener threshold $T \sim 1/(2\pi\Delta^2)$ and break
this bound by $5000\times$ via a counter-diabatic (STA) drive ---
a speedup not previously applied to QCA clocking.  (iii) The MV3
parallel-voter majority gate tolerates a stuck-at-zero defect with
$8/8$ fidelity and yields a distance-3 surface code with threshold
$p_{\rm th} \approx 10^{-2}$.  (iv) A logical Bell pair is
prepared at fidelity $1.0000$.

We then extend the platform to open quantum systems and to
programmable many-body simulation.  (v) Ohmic dephasing gives
$T_2 = 1/(2\pi\alpha k T)$, setting the clock-versus-coherence
budget for molecular QCA.  (vi) A 2D QCA lattice with patterned
couplings realizes Kitaev's toric code (4-fold ground-state
degeneracy at $L=2,3$) and the Kitaev honeycomb (4-sector
topological degeneracy decaying exponentially in $L$, confirming
non-Abelian Ising anyons).  (vii) The same lattice is a
programmable Fermi--Hubbard simulator, benchmarked against
Lieb--Wu (1D, $U/t=4$) with clean Mott physics in the doublon
density.  (viii) A QCA chain driven by random unitaries and
projective measurements undergoes a measurement-induced phase
transition at $p_c \approx 0.16$.

We validate on IBM Quantum hardware: the $L=2$ toric code ground
state on \texttt{ibm\_marrakesh} (156-qubit Eagle r3) gives
$\langle B_p \rangle = 0.868$ ($72\sigma$ above noise),
$\langle A_v \rangle = 0.643$ ($37\sigma$), and minimum
stabilizer $0.5552$.  Prior work established the quantum nature
of QCA and VQE of QCA wires on hardware~\cite{jap2026vqe}; this
work extends that foundation to a full fault-tolerant stack and
to programmable quantum simulation.
\end{abstract}
